\chapter{Conclusion}

ProblySUS stands as a demonstration of how a well-designed, multi-signal approach can meaningfully improve upon the fragmented and insufficient URL safety tools available to everyday internet users today. By combining domain intelligence, DNS analysis, URL pattern recognition, live content inspection, threat feed lookups, headless browser simulation, network profiling, privacy auditing, and tracker identification into a single unified pipeline, the system delivers a comprehensive and explainable risk assessment that no single-signal tool can match.

A key focus throughout development was building a strictly modular architecture, with each analysis concern isolated in its own independent module. This design ensures ease of testing, maintenance, and future extension without introducing cascading side effects. The frontend was developed using React~19 and Tailwind CSS~v4, focusing on clarity and immediate interpretability of results for non-technical users. The backend, implemented using Python and Flask, leverages concurrent execution and intelligent caching to deliver analysis results within a practical time frame. The browser extension further enhances usability by integrating the scanning capability directly into the user’s browsing context, allowing one-click analysis of the active tab without leaving the browser. Comprehensive testing, including unit, integration, performance, and usability evaluations, validated the reliability, accuracy, and resilience of all system components.

In summary, ProblySUS empowers users with transparent and actionable insights about the websites they visit, bridging the gap between passive blacklist-based tools and proactive, multi-dimensional threat analysis. The system establishes a strong foundation for a smarter and more accessible approach to online security in an era where scam and phishing attacks continue to grow in both volume and sophistication.